\documentclass[10pt,a4paper]{amsart}
\usepackage[margin=19mm]{geometry}
\usepackage{amsmath,amssymb,mathtools}
\usepackage[scr=rsfs,cal=euler]{mathalfa}
\usepackage{xfrac}
\usepackage[unicode,hidelinks,bookmarks=false]{hyperref}
\newcommand{\ie}{i.e.\ }
\newcommand{\cf}{cf.\ }
\newcommand{\resp}{respectively\ }
\newcommand{\Subsection}{\subsection}
\providecommand{\nobreakdash}{\nobreak\hbox{-}}
\setlength{\emergencystretch}{2em}
\multlinegap0pt
\usepackage{enumerate}
\usepackage{amssymb}
\usepackage{xcolor}
\newtheorem{lemma}{Lemma}
        \newcommand{\lrcn}{\mathsf{lrcn}}
\title{Topics, Non-Uniform Substitutions, and Variable Sharing}
\author{Shawn Standefer, Shay Allen Logan, Thomas Macaulay Ferguson}
\date{Source: arXiv:2409.08942v1 (CC BY 4.0)}
\begin{document}
\maketitle

\noindent\emph{Source and licence.} Topics, Non-Uniform Substitutions, and Variable Sharing. Version arXiv:2409.08942v1 (CC BY 4.0), distributed by arXiv under the Creative Commons Attribution 4.0 International licence, http://creativecommons.org/licenses/by/4.0/, as recorded in the arXiv licence metadata for this exact version and shown on the abstract page https://arxiv.org/abs/2409.08942v1. Full licence notice: \texttt{licenses/arxiv-2409.08942-CC-BY-4.0.txt}.\par\smallskip

\noindent\textit{Editorial scope.} Counted unit: the article's lemma Lemma:h-definition (source0.tex lines 678--689). The article's complete original proof of it is delivered with it (source0.tex lines 690--743, one \texttt{proof} environment of 54 lines). No mathematics is written, repaired or re-derived: the statement and the proof are the article's own lines, in the article's own notation, and no retained line is substituted or re-flowed.  No further source-local context is needed: every cross-reference of every retained line resolves inside the two delivered spans.  Nothing was omitted from any retained span, so the package carries no omission record.  No retained line contains a citation, so the wrapper carries no bibliography and no bibliography entry.  No retained line contains a figure environment, an image-inclusion command or drawing code, so no image is delivered and the package declares no asset dependency.  Editorial formatting only, in addition: an AMS \texttt{amsart} preamble that restates the article's own declarations and macros used by the retained lines, namely the environments \texttt{lemma} on separate counters, together with the standard packages those lines need.  The retained environments print in this document as Lemma 1 in order of appearance, whereas the article prints the same items with the numbering of its own full document.  The complete native source of the article as distributed in the arXiv e-print for this exact version is bundled unchanged as \texttt{sources/165220/source0.tex}.
\begin{lemma}\label{Lemma:h-definition}
    Suppose that $A$ and $B$ share no atoms. Define the assignment $h$ by
    \[
        h(\overline{x},p)=\left\{
        \begin{array}{rl}
            f^+_A(\overline{x}, p) & \text{ if } p \text{ occurs in $A$ and $\overline{x}=\lrcn((A\to B)[p])$} \\
            f^-_B(\overline{x}, p) & \text{ if } p \text{ occurs in $B$ and $\overline{x}=\lrcn((A\to B)[p])$} \\
            1 & \text{ otherwise}.
        \end{array}\right.
    \]
    Then if $C$ is a subformula of $A$ then $h(\lrcn((A\to B)[C]),C)=f^+_A(\lrcn((A\to B)[C]),C)$ and if $C$ is a subformula of $B$ then $h(\lrcn((A\to B)[C]),C)=f^-_B(\lrcn((A\to B)[C]),C)$
\end{lemma}

\begin{proof}
    By induction on $C$, separately for each conclusion. For atoms the result is immediate from the definition of $h$. We sample a selection of the remaining clauses and leave the rest to the reader.

    Suppose $C=D_1\land D_2$ is a subformula of $A$, and say $\lrcn((A\to B)[C])=\overline{x}c$. As $\lrcn((A\to B)[C])=\lrcn((A\to B)[D_i])$,
    it follows that
    \[
    h(\overline{x}c,C)=\min(h(\overline{x}c,D_1), h(\overline{x}c,D_2)).
    \]
    
     By IH, 
    \[
    \min(h(\overline{x}c,D_1), h(\overline{x}c,D_2))=\min(f^+_A(\overline{x}c,D_1),f^+_A(\overline{x}c,D_2))=f^+_A(\overline{x}c,C),
    \]
    which was to be proved.

    Suppose $C=D_1\lor D_2$ is a subformula of $B$, and say $\lrcn((A\to B)[C])=\overline{x}c$. Then, as $\lrcn((A\to B)[C])=\lrcn((A\to B)[D_i])$,
    \[
    h(\overline{x}c,C)=\max(h(\overline{x}c,D_1), h(\overline{x}c,D_2)). 
    \]
    
    By IH,
    \[
    \max(h(\overline{x}c,D_1), h(\overline{x}c,D_2))=\max(f^-_B(\overline{x}c,D_1),f^-_B(\overline{x}c,D_2))=f^-_B(\overline{x}c,C),
    \]
    which completes the case.

    Suppose $C=\neg D$ is a subformula of $A$, and say $\lrcn((A\to B)[C])=\overline{x}c$. As $n\overline{x}c=\lrcn((A\to B)[D])$, it follows that
    \begin{center}
    \begin{tabular}{ccc}
    $h(\overline{x}c,C)$ & = &  $1-h(n\overline{x}c,D)$\\ 
    & =& $1-f^+_A(n\overline{x}c,D)$ \\ 
    & =& $f^+_A(\overline{x}c,C)$ \\ 
    \end{tabular}
    \end{center}
    The transition from the first to the second line is justified by IH, and the remainder are justified by the definition of assignments. 
    
    Suppose $C=D_1\to D_2$ is a subformula of $B$, and say $\lrcn((A\to B)[C])=\overline{x}c$. 
    As $l\overline{x}c=\lrcn((A\to B)[D_1])$ and $r\overline{x}c=\lrcn((A\to B)[D_2])$, it follows that 
    \[
    h(\overline{x}c,C)=\max(1-h(l\overline{x}c,D_1), h(r\overline{x}c,D_2)).
    \]
    Therefore, 
    \[
    \max(1-h(l\overline{x}c,D_1), h(r\overline{x}c,D_2))=\max(1-h(\lrcn((A\to B)[D_1]),D_1), h(\lrcn((A\to B)[D_2]),D_2)).
    \]
    By IH, this implies that the right-hand side is identical to
    \[
    \max(1-f^-_B(\lrcn((A\to B)[D_1]),D_1), f^-_B(\lrcn((A\to B)[D_2]),D_2)), 
    \]
    
    which in turn is identical to $f^-_B(\lrcn((A\to B)[C]),C)$,
    as desired.

\end{proof}

\end{document}
